\section{Deep Materiality}

Tomberg remarks in his \emph{Meditations} that Christianity places the locus of sin, not in the body or material existence (such as he thinks is implied to a degree by Eastern metaphysic), but in the soul. The body, then, as such, becomes rather a focus of redemption and revelation, rather than something which needs to be shed (along with the sense of ``self''), in order to attain Enlightenment. Tomberg even refers to the personality as the ``lost sheep'' or lonely and vulnerable flock which the Lord came intending to seek and to save. These views can be linked with his comments on hermeticism differing even from ``saintliness'', in which he insists that the hermetic is like Jacob wrestling with the angel, in that it refuses to sacrifice or suborn its intellect. Instead, it wishes the intellect to be ``re-made'' but kept, in the image of holiness. This defines Tomberg's path over and against those of magical bent (such as Franz Bardon\footnote{\url{http://www.asiya.org/athenaeum/Bardon\_Magical\_Evocation.pdf}}) who might seem to believe that no essential connection exists between purification/illumination/theosis and the awakening of the vital bodies, and those who utterly reject Tomberg's aspiration to pursue something which is essentially different than the saintly path of (say) someone like Saint Cuthbert, who waded into icy sea water\footnote{\url{http://www.mainlesson.com/display.php?author=macgregor&book=christopher&story=cuthbert}} and sang psalms all night long to develop their will. Tomberg (I think) would say that the saints truncated a part of their personality, and that the magicians sacrificed a piece of their soul.

Rather than looking at this as a Hegelian synthesis, perhaps we should see Tomberg as essentially re-uniting something which was once indivisible, on a higher plane. I will call this the path of the ``hermit'', using Tomberg's own image. This (I would argue) should be seen as the path of the Middle Ages, which abrogated the reign of the saints (from St. Martin all the way to St. Guthlac) in seeking a more ``human'' path, somewhere in between ``beast'' and ``god''. For instance, Saint Martin might have helped save the Empire by engaging in martial leadership, but chose instead to pursue the ``saintly path'' of abnegation. Anyone who has read Evola would have to feel a twinge of regret when reading the life of the this saint\footnote{\url{https://en.wikipedia.org/wiki/Martin\_of\_Tours}}.

How is this path to be sought? I will not necessarily argue that it has to take place within the ``Church'' (the body of Christ) -- for a number of reasons, this is impracticable. Tomberg gives us another clue in \emph{Meditations} -- those who are ``humble'' have certainly been in the presence of God. But what, he asks, if they have not consciously known this? What then? Is the intuition invalid? No, they have experienced it somewhere else, either in a dream, or somewhere even more dreamlike.

And this leads us back to ``deep materiality''. In the ``Secret Life of Doctor Honigberger'', Eliade\footnote{\url{http://www.plural-magazine.com/article\_the-secret-of-dr-honigberger.html}} suggests that the Divine Absolute keeps most of us from developing our spiritual powers in order to both protect us, and prepare us, for that world. A certain detachment, and asceticism is necessary, because to suddenly enter it would be a sensory overload which would lead to destruction at worst, confusion and loss at best. The Middle Ages embodied the most ``natural'' and ``sane'' path because it balanced an awareness of this life as an ``antechamber'', while giving full play to the grace inherent in the innocent body. It was the soul which needed (above all) purification, not the body, or the spirit. This was the secret genius and angel which united Europe and created Christendom. As the picture reminds us, the blessed Theotokos actually nursed the man-God, who did not scorn to become a child. Let us, then, not scorn the ``days of small things'' when we are as children, and are not as we shall become. Not yet.


\flrightit{Posted on 2012-01-21 by Logres}

\begin{center}* * *\end{center}

\begin{footnotesize}
\begin{sffamily}
\texttt{Boreas on 2012-01-21 at 12:13 said:}

``For we wrestle not against flesh and blood, but against principalities, against powers, against the rulers of the darkness of this world, against spiritual wickedness in high places.'' -- Eph. 6:12

I'd add that we wrestle them in ourselves, in our personality, and when we have won ourselves and our own petty desires, the ``principalities and powers'', ``the archons of this world'' and ``the prince of this world'' reveal themselves as pure divine powers, beyond good and evil. Before this they manifest as tempters, accusers and opponents, who test the purity of our motives and intentions, our following of the divine will -- just like Satan did to Job (and later on to Christ).

\hfill

\texttt{Logres on 2012-01-21 at 17:03 said:}

That would be consistent with pursuing a ``middle path'' between abjuration pure and simple (the saints) \& premature exploration (most currents of Magick). This seems to sum up Tomberg's ``project'' in his Meditations on the Tarot. Thank you for the comment. In the end, we would look back \& see that the dark powers existed primarily because of the turning away from God to begin with (``all them that hate me, love death''). It is also consistent with the vision of Recapitulation, in which the last enemy that is conquered is Death (and its attendants).

\hfill

\texttt{Logres on 2012-01-21 at 17:08 said:}

The only warning I would add to Boreas' statement is that (historically) we have often seen how seemingly pursuing a synthetic project ends up destroying both thesis and antithesis. For example, to the uninitiate, Hegel combines classical philosophy with Christian revelation in such a way as to annihilate both (the ``Left'', Marxism). Of course, to the initiate, Hegel can be revealed to be engaged in hunting much bigger game, and in not doing what his ``defenders'' claim at all. Simone Weil also cautioned against this kind of eclecticism. It is something to consider for any ``man of the West'' in his interests which lie ``outside of'' Western tradition. There are spiritual energies at work here, and we see where most multicultural or New Age thought ends up. Hence, Gornahoor's emphasis on pursuing a ``valid'' spiritual tradition of one's very own, knowing it very well, being grounded in it. If Plato's ``Ideas'' only exist as spiritual vortex and energies generated by higher Beings, then ``playing'' with Ideas might involve great danger. This must have been the basis of picking a guide, or teacher, and of adhering to what had been ``revealed'' already: after all, why should the Absolute and the Infinite give us better guides when we won't listen to the ones we've already been given? And wasn't this the basis of medieval ``tradition''? Dante?

\hfill

\texttt{Charlotte Cowell on 2012-01-22 at 09:53 said:}

I've decided most of the answers are in Lord of the Rings :-) (having just watched the entire set of DVDs again!). Can't beat the book tho, shame the font is so small... ..

\hfill

\texttt{Charlotte Cowell on 2012-01-22 at 09:58 said:}

Also, more on topic, recently I read something by Judith von Halle about the spear wound, I think i shared this -- touches on purification/reforming of hte bodies.

\url{http://alchemical-weddings.com/alchemical-weddings/mystery-of-the-spear}


\hfill

\texttt{Logres on 2012-01-22 at 16:29 said:}

Charlotte, thanks for linking that... I was wondering what you meant by that comment.

\hfill

\texttt{Charlotte Cowell on 2012-01-22 at 17:10 said:}

oh let me see, well in my mind we are working on the formation of sustainable resurrection bodies and I felt this was relevant -- so while the saint goes the way of purification through self denial, the hermeticist wrestles with the angel. But both have the Resurrection body as an end goal, or part of it at least, do you not think that is so? I am thinking perhaps of the aspect of practice that leads to results tangible on all planes, the manifestation you spoke of recently.

\hfill

\texttt{Charlotte Cowell on 2012-01-22 at 17:12 said:}

which is also what brought me onto Lazarus, as for a while I felt it was important to think of this mystery, though I would not necessarily have been able to say why -- more it was because of the direction that we seemed to be going in with the stream of posts and comments, of consciousness

\hfill

\texttt{Charlotte Cowell on 2012-01-22 at 17:16 said:}

also there was an internal prompt to read this book of VT's, which like Christ and Sophia addresses and helps to shed light on some of the very dense and complex situations and challenges raised by MoTT

\hfill

\texttt{Charlotte Cowell on 2012-01-22 at 17:26 said:}

the being `remade' corresponds to an extract here:

\url{http://alchemical-weddings.com/alchemical-weddings/lazurus-come-forth}


The resurrection within Christianity of the Hindu and Buddhist spiritual life, to which the church owes the arising of the whole monastic movement and the founding of religious orders in late antiquity, as not the last event of its kind in churh history.

Others followed according to the law that all truth and love of the past that have timeless values are called back out of the realms of forgetting, sleep and death into the daylight of Christian spiritual life through the call that from age to age reminds, rouses and awakens.

Through this call sounding forth from Him who is the Resurrection and the Life, saying `Lazarus, come forth!' -- the most noble and valuable aspects of pagan antiquity were also resurrected.

The Platonic and Aristotelian treasury of thought arose radiant in transfigured form and inspired great spirits of the church to take up the philosophia perennis, in which lay the task of lifting up the chalice of pure human thinking and sacrificial offering to divine revelation.

For this was the essential aim of the scholastics: the raising up of the chalice of crystal clear human thinking upon the altar of Godhead -- the Godhead manifesting in divine revelation... ..

... ..for the history of Christianity is not that of adjusting to world progress -- it is the story of the return of the Prodigal Son, ie, humanity, to the Father's House. its path is not so much a forming of the future as a resurrection of the past -- not indeed the mere return of the past, but the reawakening of its eternal core of truth in all its greatness and nobility.

As the fall of the human being into sin was a descent from the heights of consciousness of the everlasting Presence of God, so does the way of reascent lead through the same stages by which the descent was made -- that is, through the stages of resurrection of the past.

And the future of the world, its end, the Last Day, is not the final triumph of the progress of natural evolution, but the state of the resurrection of the whole of the past. The Last Day is no specific `day' of the last year of the world, but is the ever present process of resurrection that culminates in the resurrection of the eternal `first day of creation'.

It is the primal creative Word: `Let there be light!' that forms the essence of the Last Judgement. yes, the world passes away -- and simultaneously it is resurrected. The eternal creative Word: `Let there be light', is the Alpha and Omega, the first and the last, and the eternally present (Revelation xxi, 6).

\hfill

\texttt{Logres on 2012-01-22 at 20:42 said:}

ntique and ``difficult'' transformations.


\hfill

\end{sffamily}
\end{footnotesize}

